\section{理论框架}\label{sec:theory}

\subsection{外部性与公共品}

燃放是消费产生的负外部性：燃放者获得效用，成本由同城居民在数小时内共同承担。清洁空气是非排他、非竞争的，因此即便不存在燃放，其供给也会不足。这两点解释了政府为什么干预，但不解释政府\textbf{用什么样的工具}干预。后者才是本文的问题。

\subsection{价格与数量：Weitzman 条件}

按 Weitzman（1974），在边际减排成本不确定的条件下，数量工具优于价格工具的条件是\textbf{边际损害曲线比边际减排成本曲线更陡}。

燃放的边际损害高度依赖时间（深夜、跨零点）、地点（密集居民区）与气象（静稳天气）。这种依赖使损害曲线在时间上近乎不连续：一年 365 天，损害集中在十几个小时。相比之下，边际减排成本曲线——也就是“少放一箱”的机会成本——在时间上的变化要平缓得多。

\textbf{因此本文的读法是：这是一个数量工具占优的极端情形。} 地方政府选择禁令，从这个标准看是理性的，而不是无知或因循。

但这里有一处必须谨慎。燃放的需求刚性强（习俗、一次性、单价低），这意味着边际减排成本曲线本身也可能很陡。\textbf{两条曲线谁更陡，本文不将其作为结论，只作为讨论。} 本文能确立的是更弱但更可靠的一点：损害在时间与空间两个维度上的变异，远大于任何全国统一、按年适用的工具所能表达的变异。

\subsection{解析度：一个可操作的定义}\label{sec:resolution}

上面这句话在文献里常以比喻出现。本文把它写成一个可以直接计算的对象。

设损害（或本文实际测量的尖峰增量）$D(c,t)$ 定义在城市 $c$ 与年份 $t$ 上，总方差为 $\sigma^2$。任一政策工具都会在 $(c,t)$ 空间上诱导一个\textbf{划分} $\mathcal{P}$：同一个格子内的所有 $(c,t)$ 只能被施加同一个政策参数——同一个税率、同一个禁令强度。定义该工具的\textbf{解析度损失}为

\begin{equation}
L(\mathcal{P}) \;=\; \mathbb{E}\big[\operatorname{Var}\big(D \mid \text{cell}\big)\big]
\;=\; \sigma^2 \;-\; \operatorname{Var}\big(\mathbb{E}[D \mid \text{cell}]\big).
\label{eq:loss}
\end{equation}

$L$ 就是“工具够不到的那部分变异”。它有三个直接的性质：

\begin{enumerate}
  \item \textbf{全国统一、按年适用的从价税}只有一个格子，$\mathbb{E}[D\mid\text{cell}]$ 是常数，故 $L$ 等于总方差——\textbf{工具对损害的变异完全没有分辨率}。
  \item \textbf{把税率按城市设定}（仍由中央统一）得到 $|C|$ 个格子，$L$ 降为总方差减去城市间方差。
  \item \textbf{按城市 $\times$ 年份设定}得到 $|C|\times|T|$ 个格子，$L=0$（在年度频率上）。市级禁放/限放之所以更接近这一端，是因为它同时划定地点与时段。
\end{enumerate}

这个定义有三处必须说明的边界。第一，$D$ 是\textbf{已实现}的损害，其中包含气象等不可预测成分；工具能利用的只是其中可预期的部分，因此 $L$ 是一个\textbf{上界意义上的}分辨率缺口，不是可实现的福利损失。第二，禁令改变的不是价格而是可行集，把“禁令强度”映射为格子内的一个常数是一种简化。第三，式~\eqref{eq:loss} 只刻画“能否对准”，不刻画对准之后的效率——后者仍取决于 Weitzman 意义上的斜率比较。

尽管有这三处边界，$L$ 的价值在于：\textbf{它把一个关于“工具形式”的定性判断，变成了可以用一份面板数据直接算出来的三个百分比。} 第~\ref{sec:magnitude}~节将给出这个计算。

\subsection{工具成本的可感知性}\label{sec:visibility}

第三块理论来自一个已有文献家族。Beiser-McGrath, Bernauer \& Prakash（2023）在北京与新德里做的联合实验发现，命令控制型工具之所以受欢迎，正因为它\textbf{把政策成本模糊化}（obfuscate policy costs），而市场型工具\textbf{强加可见的成本}（impose visible costs）；提高燃油税会降低支持，但把收入用于补贴电动车或公交可以消除这一负面影响。Harrington, Morgenstern \& Sterner（2004）把“成本显现”（cost revelation）列为比较工具的标准之一。Buchanan \& Tullock（1975）从公共选择角度给出同一方向：在位的污染者偏好配额——它有租金、有进入壁垒——而反对税。Oates \& Baumol（1975）则指出一种不对称：设定配额时可以“不负责任地严格”而不注意其沉重成本，但要把税率提到足以实现同样目标的水平，很难不被注意到。

把这个框架搬到本文的工具对上，得到一个可检验的制度判断：\textbf{禁令的成本落在巡查、蹲守与查处的警力工时里，不进入任何价格，也不被单独核算；而加税的成本以涨价的形式直接落在每个消费者身上。} 因此政府在选择工具时比较的，可能不是成本有效性，而是\textbf{成本的政治可见性}。

\begin{quote}
\textbf{定位说明}：这个概念不是本文发明的。本文的贡献是把这一已有框架应用到一个新的工具对（全国从价的消费税 vs.\ 市级禁放限放）与一个新的场景（中国的烟花爆竹），并用政策文件与制度证据说明禁令的成本形态。
\end{quote}

\subsection{记账规则}

本文比较成本时遵循两条规则：\textbf{警力、人力、设备是真实资源成本}；\textbf{罚款与税收是转移支付，不计入社会净成本}。销售额的变化也不等于社会福利的变化。这两条会直接影响第~\ref{sec:comparison}~节的读法。
